% \clearpage


\section{Related Work}

\paragraph{Video Generative Models.}
Recent advancements in video generative models have mainly been driven by models that leverage latent representations derived through dimensionality reduction. These representations can be either continuous~\cite{kingma14,nvidia2024cosmos} or quantized~\cite{oord17,esser21}, and the success of video generation heavily depends on how efficiently these representations capture spatiotemporal information. One of the key challenges lies in ensuring that the mappings to and from latent space preserve both visual fidelity and semantic coherence. While scaling to large datasets improves generalization~\cite{agarwal2025cosmos}, challenges remain in preserving both visual realism and temporal consistency.

Autoregressive and masked generative models using quantized tokens have demonstrated strong modeling of temporal dependencies and motion dynamics~\cite{yan2021videogpt,moing21,ge2022long,seo22,hawthorne22,micheli2022transformers,teco22yan,villegas22phenaki,yu2023magvit}. However, their sequential generation processes lead to slow inference and error accumulation, which can limit their scalability for long sequences or complex multi-agent scenes. In contrast, latent diffusion models offer a more parallelizable alternative, using iterative denoising to generate full video sequences. These models, including Stable Video Diffusion~\cite{blattmann2023stable}, DIAMOND~\cite{alonso2024diffusionworldmodelingvisual}, Sora~\cite{openai2024sora}, MovieGen~\cite{polyak2024movie}, Gen-3~\cite{runway2024gen3}, and Cosmos~\cite{agarwal2025cosmos}, have achieved impressive quality on text-to-video and image-to-video tasks, especially in open-domain settings. 

However, most of these works prioritize visual aesthetics and domain-agnostic generation, offering limited support for structured control over scene elements or agent behavior. This makes them less suitable for applications (such as autonomous driving) that require precise, semantic, and multi-modal scene control across space and time. In this context, GAIA-2 aims to bridge the gap between high-fidelity generation and domain-specific controllability required for safety-critical systems.


\paragraph{Generative World Models in Autonomous Driving.}
Unlike general-purpose video generation, generative world models for autonomous driving must satisfy stricter requirements: multi-agent interaction, ego-motion control, environmental diversity, and multi-camera coherence. Early work in this space includes GAIA-1~\cite{gaia1-2023}, which introduced text and ego action conditioning in a discrete world model with a video diffusion decoder to improve temporal consistency. CommaVQ~\cite{comma2023vq} similarly leveraged a causal transformer on discrete tokens for ego-motion control. Although these methods took important steps toward controllability, they often targeted single-camera settings and partial controllability via action or text conditioning, limiting their applicability to comprehensive driving simulations.

Subsequent approaches adopted latent diffusion for higher fidelity generation and introduced more flexible control. DriveDreamer~\cite{wang2023drivedreamer} adopts a latent diffusion model conditioned on 3D bounding boxes, HD maps, and ego actions, and further incorporates an action decoder to predict future ego actions. Drive-WM~\cite{wang2024driving} extends this approach to a multi-camera setting (up to six cameras surrounding the ego vehicle) and introduces controllability over ego actions, dynamic agents, and environmental conditions, such as weather and lighting. UniMLVG~\cite{chen2024unimlvg} similarly enables multi-view generation conditioned on text, camera parameters, 3D bounding boxes, and HD maps, while MaskGWM~\cite{ni2025maskgwm} extends UniMLVG to support longer video generations. Vista~\cite{gao2024vista} employs latent diffusion for generating high-resolution, long-duration videos, and Delphi~\cite{ma2024unleashing} uses a data-driven approach to guide video generation toward failure cases in driving models, thereby improving the synthetic data's training utility. More recently, GEM~\cite{hassan2024gem} further generalizes scene control to include ego motion, object dynamics, and even human poses, suggesting broader applicability across domains such as human motion synthesis and drone navigation. A recent line of work has explored 4D geometry-aware generation. DriveDreamer4D~\cite{zhao2024drivedreamer4d} uses video generative models to synthesize videos along novel trajectories and combines real and generated footage to optimize a 4D Gaussian Splatting model for closed-loop simulation. Similarly, DreamDrive~\cite{mao2024dreamdrive} leverages generative priors to construct high-quality 4D scenes from in-the-wild driving data.

Despite these advancements, most of these solutions only address subsets of the requirements for realistic driving simulation. Some are limited to single-camera settings, others lack structured agent-level control, and few integrate multi-view generation and agent-level semantics into a single framework. Additionally, many do not support fine-grained editing tasks such as inpainting or targeted scene modification, essential for data augmentation and scenario variation in simulation loops.

Motivated by these limitations, \textbf{GAIA-2} is designed as a unified latent diffusion framework that brings together: 
\begin{itemize}
    \item multi-camera generation (up to five views); 
    \item structured and semantic conditioning over ego action, dynamic agents, and scene-level metadata; 
    \item flexible inference modes, including generation from scratch, inpainting, scene editing, and long-horizon rollouts; and 
    \item support for external latent spaces like CLIP and driving scenario embeddings.
\end{itemize} 
By combining continuous latent representations with large-scale training on a geographically and environmentally diverse dataset, GAIA-2 delivers both visual fidelity and scenario controllability. It fills a critical gap between general-purpose diffusion models and the domain-specific needs of autonomous driving, namely, the ability to simulate realistic, controllable, multi-view driving scenes for robust, scalable training and evaluation.
